% ============================================================
% ch12_mio_observatory_results.tex
% MIO Observatory Results
% Part of: Tetrad-Based Departure Decomposition (Jiwon Park)
% W8D5 draft — 2026-04-19 (MANU-CH12-NEW first landing)
% Sections shipped in this landing: 12.0, 12.2, 12.6, 12.7.
% Sections 12.1, 12.3, 12.4, 12.5, 12.8 remain blocked on
% bass_py / HTT outputs (plan v1.2 §13.2 "후행") and will be
% drafted in a later week.
% ============================================================

\chapter{MIO Observatory Results}
\label{ch:mio-observatory}

This chapter reports the output of the Model-Independent
Observatory (MIO), the fourth pillar of the solver ecosystem
introduced in Chapter~\ref{ch:pipeline}.  MIO differs from the
forward solver (Chapter~\ref{ch:results}), the admissibility
diagnostics
(Chapter~\ref{ch:bianchi-bounds}--\ref{ch:robustness}), and the
model-dependent inference
(Chapter~\ref{ch:results}--\ref{ch:discussion}) in a single
structural way: it reports departure variables extracted from
data \emph{without assuming any specific cosmological model},
and it never returns a posterior sample, an evidence
contribution, or a truth certificate.

The chapter is organised around the five non-parametric
reductions identified in the research plan (v3 §11.14.1): shear
extraction (§12.1), directional coherence (§12.2), FLRW-null
tension and $x_C$ estimation (§12.3), evidence anatomy (§12.4),
predictive residuals (§12.5), HTT cross-validation (§12.6),
scope and limits (§12.7), and a summary comparing MIO against
HTT outputs (§12.8).  Sections §12.1, §12.3, §12.4, §12.5, and
§12.8 depend on forward-solver outputs ($K_\ell$ atlas, BiPoSH
coefficients, and HTT posteriors) and are deferred to a later
landing; the present chapter opens with the four sections that
are independent of those dependencies.

% ============================================================
\section{Philosophy: a Model-Independent Observatory}
\label{sec:mio-philosophy}
% ============================================================

MIO is the Model-Independent Observatory.  The name is
deliberate and load-bearing: MIO is neither a solver (it does
not evolve a physical system forward in time), nor an
inference engine (it does not compute posteriors or
log-evidences), nor a gatekeeper (it does not issue truth
certificates).  MIO is a \emph{physics-producing observatory}
in the same functional category as a spectrograph or a
CMB-temperature calibrator: it emits a report that describes
\emph{what is directly visible in the data, without assuming a
model}, and leaves the model-dependent interpretation to the
inference layer.

\paragraph{The single deliverable.}

Every MIO reduction produces one object of a single type, the
\texttt{MioCertificate} (Chapter~\ref{ch:error-hierarchy},
Section~\ref{sec:err-mio-semantic}).  The object is a frozen
dataclass: it carries a \texttt{report\_type}, a
\texttt{probe\_name}, a \texttt{channel}, a
\texttt{departure\_variables} dictionary of model-independent
measurements, a \texttt{reduction\_status} three-valued
enumeration that records whether the reduction is
theory-direct, theory-approximate, or diagnostic-only, and a
provenance block.  The certificate has no field with
posterior semantics, and its schema is locked by a hash digest
so any later drift fails loudly in CI.

\paragraph{Three structural commitments.}

Three architectural commitments distinguish MIO from the
surrounding code:

\begin{enumerate}
\item \emph{No model assumption.}  Each MIO reduction computes
  its departure variable with no reference to any specific
  Bianchi type or parameterised model.  The $\Sigma^2_{\rm
  MIO}(\ell)$ extraction of §12.1 relies only on the
  observed $C_\ell$ spectrum, a $\Lambda$CDM baseline, and a
  forward-solver kernel $K_\ell$; the directional coherence
  of §12.2 relies only on the sky positions and
  pointing-cone widths of five independent probes; the FLRW
  tension metric of §12.3 relies only on FLRW-null mock
  draws and a test statistic.  None of these reductions can
  be coerced into preferring one Bianchi type over another.
\item \emph{No posterior output.}  The
  \texttt{MioCertificate} schema does not contain a field
  named \texttt{posterior}, \texttt{posterior\_samples},
  \texttt{log\_likelihood}, or any synonym.  The declared
  method
  \texttt{MioCertificate.as\_posterior\_bundle()} raises
  \texttt{NotImplementedError} with a pointer to the HTT
  posterior interface, so any attempt to convert a
  certificate into a posterior fails loudly.
\item \emph{No master score.}  Outputs from MIO, HTT, and
  TSC never combine into a single merged scalar
  (Chapter~\ref{ch:error-hierarchy},
  Section~\ref{sec:err-g19}).  Cross-checks are permitted
  and encouraged; merges are architecturally forbidden.
\end{enumerate}

\paragraph{What MIO is not.}

It is easier to describe MIO by exclusion.  MIO is
\emph{not} a central adequacy module that attests on behalf
of the pipeline that a model is correct.  It is \emph{not} a
post-hoc re-weighting layer that adjusts HTT evidence by
MIO-derived scores.  It is \emph{not} a truth certificate,
in the sense that no successful construction of a
\texttt{MioCertificate} is a claim that any particular
cosmological model has been validated.  These exclusions
correspond to three load-bearing defences at the code level:
a frozen schema that forbids posterior fields, a
generator-site keyword scan that refuses any input named
after a posterior quantity, and an ingestion-site type-check
on HTT entry points that refuses any
\texttt{MioCertificate} instance.  The exclusions are not
rhetorical; they are enforced tests in the CI battery.

\paragraph{Why a separate chapter.}

MIO deserves a separate chapter, rather than a subsection of
Chapter~\ref{ch:pipeline} or Chapter~\ref{ch:robustness},
for two reasons.  First, its outputs are a new class of
scientific object in the thesis: model-independent departure
reports that complement the model-dependent evidence
comparisons of Chapter~\ref{ch:results}.  Second, the
architectural constraints enforced on MIO
(Chapter~\ref{ch:error-hierarchy},
Sections~\ref{sec:err-mio-semantic}--\ref{sec:err-epistemic})
require that MIO output be presented alongside a clear
declaration of its scope and its non-overlap with HTT; a
separate chapter makes that declaration unambiguous to the
reader.

\paragraph{Relationship to the 596-line Chapter~11 extension.}

Chapter~\ref{ch:error-hierarchy} introduced the MIO
architectural rules as extensions of the five-level error
hierarchy.  The present chapter applies those rules to the
three non-parametric reductions whose results are already
in hand: directional coherence (§12.2, HJ-02a data-only),
HTT cross-check protocol (§12.6, TSC-06 landed), and the
re-statement of scope (§12.7).  The remaining five sections
will populate Chapter~\ref{ch:mio-observatory} as the
upstream dependencies land; the present landing is therefore
a partial deliverable whose completeness is strictly
bounded by the Phase-J roadmap (v3 §0.4).


\paragraph{The observatory metaphor.}

The functional analogy to keep in mind throughout the
chapter is that of a physical instrument.  A CMB
spectrograph reports a temperature in kelvin; a photometer
reports a photon count; neither claims that the measured
value validates any particular cosmological model.  The
MIO reductions are built in the same spirit: each emits a
measurement of a departure variable together with the
caveats that scope its interpretation, and the task of
deciding what the measurement \emph{means} for a given
cosmological model is left to the inference layer.  The
\texttt{MioCertificate} is the digital analogue of a
FITS-formatted calibration record: structured, provenanced,
and blind to the theory prior.

\paragraph{Distinction from the TSC suite.}

The distinction from the Tangency--Stability--Chart suite
(Chapter~\ref{ch:pipeline},
Section~\ref{sec:pipeline-tsc}) is equally important.  TSC
evaluates whether a configuration of
$(\epsilon_1, \epsilon_2, \epsilon_3)$ is physically
realisable under the Paper~I $T_{\rm eff}$ chart; its output
is a Boolean admissibility flag and the three-bound /
Michaelis-Menten SSOT mirrors.  MIO, by contrast, consumes
\emph{observations} rather than
$(\epsilon_1, \epsilon_2, \epsilon_3)$ triples, and its
output is a \texttt{MioCertificate} populated with
data-derived departure variables.  TSC asks ``is this
configuration physically allowed?'' MIO asks ``what do the
data say when we do not assume a model?''  Both feed the
cross-check protocol of §12.6; neither is a substitute for
the other.


% ============================================================
\section{Cross-channel directional coherence}
\label{sec:mio-coherence}
% ============================================================

The directional-coherence reduction answers one question on
data alone: \emph{do the five independent directional probes
of cosmic anisotropy---the CMB dipole, the CatWISE
high-redshift galaxy dipole, the NVSS+RACS radio dipole, the
Cosmicflows-4++ principal-tide direction, and the Planck
BipoSH preferred axis---point to a common axis on the sky,
or are they mutually isotropic?}  The reduction is the
single MIO diagnostic that consumes neither a forward
$K_\ell$ atlas nor an HTT posterior (v3 §16.2
``완전 독립'' row; appendix~A35); it therefore lands
before the remainder of Phase~J.

\paragraph{Inputs.}

Five probe directions live in
\texttt{mio.coherence.directional.STANDARD\_PROBES} as a
single-source-of-truth table.  Each entry supplies a
galactic longitude $\ell$, a latitude $b$, a pointing-cone
width $\sigma_{\rm cone}$ (degrees), and an inverse-variance
weight.  The table is reproduced in summary form in
Table~\ref{tab:mio-probes}; the full provenance is recorded
in appendix~A35.

\begin{table}[h]
\centering
\setlength{\tabcolsep}{5pt}
\small
\begin{tabular}{lrrrl}
\hline\hline
Probe & $\ell$ (deg) & $b$ (deg) & $\sigma_{\rm cone}$ (deg)
      & Reference \\
\hline
CMB      & 264.02 & $+48.25$ & 0.5  & Planck 2018 int.\ \\
CatWISE  & 238.2  & $+28.8$  & 10.0 & Secrest+2020 \\
Radio    & 253.0  & $+25.0$  & 10.0 & Singal 2011 $^\dagger$ \\
CF4++    & 301.0  & $-30.0$  & 15.0 & Tully+2023 $^\dagger$ \\
BipoSH   & 222.0  & $-22.0$  & 20.0 & Planck 2015 $^\dagger$ \\
\hline\hline
\end{tabular}
\caption{Five directional probes consumed by
\texttt{mio.coherence.directional}.  $\dagger$: the Radio,
CF4++, and BipoSH $\sigma_{\rm cone}$ values are plan-suggested
placeholders awaiting literature citations (v3 §16.2 FM2);
the common-axis resultant is dominated by the $\sigma=0.5^\circ$
CMB direction and the placeholder-driven uncertainty is
therefore negligible for the present reduction.  Full
provenance in appendix~A35.}
\label{tab:mio-probes}
\end{table}

\paragraph{Reduction.}

Four derived statistics are computed, each model-independent:

\begin{enumerate}
\item The inverse-variance-weighted \emph{spherical mean}
  $(\bar\ell, \bar b, R)$ of the five probes, via
  \texttt{common.sky\_geometry.spherical\_mean}.  The
  resultant length $R \in [0, 1]$ satisfies $R=1$ for
  exact alignment and $R=0$ for antipodal cancellation.
\item The \emph{Monte-Carlo isotropy $p$-value}, defined as
  the fraction of isotropic mock draws whose resultant
  length exceeds the observed $R$.  The estimator uses
  Lidstone smoothing, $p_{\rm iso} = (k+1) / (N_{\rm mock}
  + 1)$, so the reported $p_{\rm iso}$ is a strict upper
  bound that never collapses to zero when no mock exceeds
  the observation.
\item The $5\times 5$ \emph{pairwise separation matrix} in
  degrees, which supplies a direct literature cross-check
  against the published $27.8^\circ$
  CMB--CatWISE separation (Secrest et al.\ 2020).
\item The \emph{common-axis $\chi^2/{\rm d.o.f.}$},
  $\chi^2 = \sum_i (\Delta_i / \sigma_i)^2$, with
  $\Delta_i$ the angular separation between probe $i$ and
  the fitted common axis.
\end{enumerate}

\paragraph{Landed artefact (2026-04-19, seed 20260419).}

Serialised to \texttt{mio\_directional\_coherence\_v1.json}
under the \texttt{mio\_} filename-prefix rule (v3 §14.1):

\begin{center}
\begin{tabular}{ll}
\hline\hline
Statistic & Value \\
\hline
$(\bar\ell, \bar b)$ & $(260.9^\circ, 42.5^\circ)$ \\
$R$ & $0.9990$ \\
$p_{\rm iso}$ & $\approx 3 \times 10^{-4}$ \\
$\chi^2 / {\rm d.o.f.}$ & $9.45$ (d.o.f.\ $= 3$) \\
Pairwise max separation & $80.4^\circ$ (Radio $\leftrightarrow$ BipoSH) \\
\hline\hline
\end{tabular}
\end{center}

The low $p_{\rm iso}$ and high $R$ reflect the informal
literature consensus that the matter-frame dipole
directions cluster toward Hydra--Centaurus--Norma.  The
large $\chi^2/{\rm d.o.f.}$ records a tension that the
common-axis fit does not absorb: CF4++ and BipoSH deviate
from the CMB axis by $60\text{--}80^\circ$, so the
``common axis'' is an inverse-variance-weighted best-fit,
not a universal alignment.  The relative weight of the
$\sigma_{\rm cone}=0.5^\circ$ CMB direction dominates
$(\bar\ell, \bar b)$ by construction.

\paragraph{MioCertificate wiring.}

The reduction serialises through
\texttt{to\_mio\_certificate(\dots)} into a
\texttt{MioCertificate} with
\texttt{report\_type="directional\_coherence"},
\texttt{channel="dipole"},
\texttt{reduction\_status="theory-direct"}, and three
placeholder caveats in \texttt{domain\_caveats}
(\texttt{radio\_sigma\_cone\_plan\_placeholder},
\texttt{cf4pp\_sigma\_cone\_plan\_placeholder},
\texttt{biposh\_sigma\_cone\_plan\_placeholder}).  The
certificate is emitted to disk under the
\texttt{mio\_directional\_coherence\_v1.json} filename and
loaded downstream by the TSC-06 advisory cross-check
(§12.6).  The certificate is not promoted to a posterior
draw at any point in the pipeline, and it is not
interpretable as a truth certificate for any specific
Bianchi type.

\paragraph{Scope-of-validity.}

Two operating assumptions are documented in A35.7 and must
accompany any downstream use of the reduction:

\begin{itemize}
\item Each probe reports its direction in the matter frame
  (heliocentric frame boosted into the local CMB rest
  frame).  Probes reporting in the geometry frame require a
  boost-correction step that
  \texttt{spherical\_mean} does not apply; the reduction
  currently trusts the published matter-frame convention
  and flags the assumption in
  \texttt{channel\_caveats}.
\item The $\sigma_{\rm cone}$ numbers assume an isotropic
  Gaussian cone around the reported direction.  Probes
  with asymmetric or non-Gaussian pointing uncertainty
  would inflate $\chi^2/{\rm d.o.f.}$; the full
  HJ-02 upgrade path (post-Phase~J) will switch to
  per-probe posterior draws when available.
\end{itemize}

\paragraph{What the reduction does not claim.}

The directional-coherence certificate claims $R=0.9990$ and
$p_{\rm iso} \approx 3 \times 10^{-4}$ on the specific
probe set above under the plan-suggested $\sigma_{\rm cone}$
placeholders.  It does \emph{not} claim that any specific
Bianchi type is preferred, that the observed alignment is
cosmological rather than selection-induced, or that the
certificate constitutes a truth certificate on the
anisotropic hypothesis.  The model-dependent interpretation
of the axis---whether Bianchi-V tilt, Bianchi-I shear, or
khronon-gradient geometry---is deferred to the evidence
comparison of Chapter~\ref{ch:results} and the HTT
cross-check of §12.6.


% ============================================================
\section{HTT $\leftrightarrow$ MIO cross-validation}
\label{sec:mio-htt-crosscheck}
% ============================================================

The cross-validation section records the protocol connecting
the Bianchi-typed Hypothesis Toolkit (HTT), the
Tangency--Stability--Chart suite (TSC), and the MIO reductions
of the previous sections.  The protocol is governed by the
\emph{hard separation} rule (G19) of
Chapter~\ref{ch:error-hierarchy},
Section~\ref{sec:err-g19}: the three adequacy domains never
merge into a single scalar, only cross-check.  This section
catalogues the wired cross-check channels, their frozen
contracts, and the four-layer enforcement that prevents a
later drift toward merging.

\paragraph{The five-property definition.}

Appendix~A34 defines a G19-compliant cross-check as any
object that simultaneously satisfies five structural
properties:

\begin{enumerate}
\item Produces two numerically comparable scalars from two
  \emph{independently implemented} code paths.
\item Does not combine the two scalars into any merged,
  summed, or averaged estimator.
\item Carries a frozen \texttt{is\_cross\_check=True} tag
  that cannot be set to \texttt{False} at construction.
\item Exposes a loud-fail guard
  (\texttt{assert\_cross\_check\_consistent}) that refuses
  to return normally on numerical mismatch or band
  violation.
\item Is consumed downstream only for agreement reporting,
  never as a likelihood term, a posterior update, or an
  evidence contribution.
\end{enumerate}

Any object missing one of the five properties is not a
cross-check; any object exposing a merged estimator is a
G19 violation and must fail the build.

\paragraph{The wired channel catalogue.}

Five G19-compliant channels are presently wired in the
codebase (appendix~A34.3):

\begin{center}
\small
\begin{tabular}{lll}
\hline\hline
Channel & Contract artefact & Status \\
\hline
TSC-06 filling fraction
  & \texttt{FFCrossCheckReport}
  & Landed W7D5 \\
TSC-03 three-bound hierarchy
  & \texttt{all\_agree} Boolean
  & Landed W7D2 \\
TSC-05 MM SSOT mirror
  & \texttt{assert\_mirror\_matches\_bass\_ssot}
  & Landed W7D3 \\
PR13AH spherical-mean
  & \texttt{ChannelSummary} diff
  & Landed W4 R1 \\
MIO HJ-02a directional
  & \texttt{MioCertificate} (advisory)
  & Landed W6 \\
\hline\hline
\end{tabular}
\end{center}

TSC-06 is the single channel that exposes a frozen
\texttt{is\_cross\_check=True} dataclass with a matched
numerical pair; the other four satisfy the five properties
through sibling mechanisms (Boolean flag, literal
bit-identity assertion, \texttt{ChannelSummary} diff
regression, advisory-only \texttt{domain\_caveats} flag).
The HJ-02a row is \emph{advisory only} because no matched
numerical partner currently exists on the HTT side; the
cross-check is consumed as a domain caveat rather than as a
scalar comparison.

\paragraph{TSC-06 as the archetype.}

The filling-fraction cross-check is the model case for the
pattern the plan expects future channels (HJ-04 evidence
anatomy, HJ-03 FLRW tension) to follow.  The
\texttt{FFCrossCheckReport} dataclass is frozen and carries
two separate fields:

\begin{verbatim}
@dataclass(frozen=True)
class FFCrossCheckReport:
    F_Bayes_tsc:        float   # Path A (tsc diagnostic)
    F_Bayes_htt_mean:   float   # Path B (htt mc_posterior)
    F_Bayes_htt_median: float
    abs_difference:     float
    rel_difference:     float
    within_published_band: bool
    is_cross_check:     bool = True
    scenario:           str
    ...
\end{verbatim}

Construction raises \texttt{ValueError} if
\texttt{is\_cross\_check} is set to \texttt{False}, and the
report exposes no merged \texttt{F\_Bayes\_combined} field.
The two code paths agree at ${\rm rtol} < 10^{-6}$ on the
published S3 scenario and both land in the Planck/CatWISE
published band $F_{\rm Bayes} \in [0.068, 0.118]$.  Loud-fail
is handled by \texttt{assert\_cross\_check\_consistent},
which raises \texttt{CrossCheckMismatch}---a direct subclass
of \texttt{AssertionError}---if the relative difference
exceeds the tolerance or the band is violated.

\paragraph{Four-layer G19 enforcement.}

A merged master score is prevented at four independent
layers (Chapter~\ref{ch:error-hierarchy},
Section~\ref{sec:err-g19}):

\begin{itemize}
\item \textbf{Type layer.}  Frozen schemas forbid the
  offending field names.  \texttt{MioCertificate} has no
  \texttt{posterior} field; \texttt{FFCrossCheckReport}
  has no \texttt{F\_Bayes\_combined} field.  Any PR that
  adds such a field breaks the schema hash and fails
  \texttt{test\_miocertificate\_schema\_frozen} or its
  sibling.
\item \textbf{Generator-site layer.} The
  \texttt{build\_mio\_certificate} factory scans its
  keyword arguments for the substring
  \texttt{"posterior"} and raises before any instance is
  built.
\item \textbf{Ingestion-site layer.} HTT likelihood entry
  points refuse any
  \texttt{MioCertificate} input at type-check time.  The
  refusal is tested by
  \texttt{test\_htt\_cannot\_ingest\_miocertificate\_as\_likelihood}.
\item \textbf{Repo-wide scan layer.} A regex-based static
  scan (\texttt{test\_g19\_enforcement}) walks the
  production subpackages for patterns such as
  \texttt{mio\_score + htt\_lnB} or
  \texttt{np.mean([mio\_score, htt\_lnB])} and fails the
  build on any unreviewed hit.
\end{itemize}

\paragraph{What a merge would break.}

Appendix~A40 catalogues four architectural failures that a
relaxation of G19 would introduce: collapse of independent
witnesses, miscounting of shared uncertainty, contamination
of HTT likelihoods, and destruction of the TSC
realisability gate.  Each is pipeline-level: the numerics
might still satisfy the five-level error budget of
Chapter~\ref{ch:error-hierarchy} while the merged object
ceases to be interpretable.  The cross-check protocol is
therefore not an aesthetic choice---it is the rule that
keeps the three adequacy reports interpretable at all.

\paragraph{Present disposition.}

All five wired channels pass in the current test battery
(994 tests on the touched surface at Week~7 boundary; 0
failures).  TSC-06 agrees at ${\rm rtol} < 10^{-6}$ on the
S3 scenario; TSC-03 agrees at ${\rm rtol} = 10^{-10}$ on
the MES coefficient table; TSC-05 agrees at bit-identity
on the Route-B constants; PR13AH spherical-mean agrees at
the W4 R1 regression anchor; and the MIO HJ-02a advisory
caveats are consumed non-numerically by the HTT dipole
likelihood.  No merged channel exists, none is planned,
and the enforcement stack refuses any addition that would
introduce one.


% ============================================================
\section{Scope and limitations}
\label{sec:mio-scope}
% ============================================================

This section re-states, in one compact place, what a MIO
output claims and what it does not.  The intent is to give
the reviewer---and any later reader inclined to treat the
\texttt{MioCertificate} as a stand-alone attestation---a
single reference for the architectural constraints.

\paragraph{MioCertificate is a diagnostic report.}

A \texttt{MioCertificate} is a departure-variable report
emitted by one of the five MIO reductions.  It is not a
posterior draw, it is not an evidence contribution, and it
is not a truth certificate.  The three exclusions are
enforced simultaneously at four code layers
(Chapter~\ref{ch:error-hierarchy}, Section~\ref{sec:err-g19}):
a frozen schema that contains no field with posterior
semantics; a generator-site scan that forbids
posterior-adjacent keyword arguments; an ingestion-site
type-check that prevents any HTT likelihood from consuming
a certificate; and a repo-wide regex scan that fails the
build on any attempt to sum a MIO scalar with an HTT
log-evidence.

\paragraph{What a certificate claims.}

A certificate claims only what its
\texttt{departure\_variables},
\texttt{adequacy\_indicators}, and
\texttt{consistency\_metrics} record.  For HJ-02a (§12.2),
the claim is that the five-probe resultant length is
$R=0.9990$ under the plan-suggested $\sigma_{\rm cone}$
placeholders and that the isotropy $p$-value is
$\approx 3 \times 10^{-4}$ at the recorded seed.  For
HJ-05a-lite (appendix~A38), the claim is a pixel-count
ratio $f_{\rm sky,\,effective}$ together with the
mask-provenance and ZoA-half-angle caveats that feed
\texttt{domain\_caveats}.  The claims are always scoped to
the probe-set, the reduction, and the caveats recorded in
the certificate; they are never scoped to a specific
cosmological model.

\paragraph{What a certificate does not claim.}

A certificate does not claim that:

\begin{itemize}
\item any specific Bianchi type is correct, preferred, or
  ruled out;
\item the observed departure is cosmological in origin
  rather than selection-induced or systematic;
\item the MIO reduction is an adequate replacement for the
  HTT posterior, the TSC admissibility Boolean, or the
  matched-complexity report;
\item a successful construction of the certificate
  constitutes an attestation that the solver ecosystem has
  validated any claim about the Universe.
\end{itemize}

The \texttt{reduction\_status} field makes one of these
bounds explicit: a reduction tagged
\texttt{diagnostic-only} is flagged as too model-near to
carry a model-independent claim on its own, and it is
retained as a diagnostic rather than promoted to a
production output.

\paragraph{What a MioCertificate is not.}

A \texttt{MioCertificate} is not a truth certificate.  It
is not a post-hoc re-weighting instrument for HTT
evidence.  It is not a merge input for the TSC
admissibility gate.  Each of these exclusions is a
load-bearing architectural commitment; none is optional.
The legacy/mio subtree preserves an earlier draft in which
some of these exclusions were absent, and the banner
attached to \texttt{legacy/README.md} (landed W5D7) warns
explicitly against importing that draft into active code
paths.

\paragraph{Where the full scope lives.}

The full scope declaration for each of the five MIO
reductions lives in appendices~A32 (schema), A34 (G19
cross-check protocol), A35 (directional coherence), A38
(masked-sky caveats), A39 (per-module epistemic
ownership), and A40 (G19 architectural stance).  The
present section summarises those appendices to a single
claim: \emph{MIO output is a data-derived departure
report, and nothing else; in particular, a
\texttt{MioCertificate} is not a truth certificate on
any cosmological hypothesis}.


\paragraph{A compact table of non-claims.}

It is useful to consolidate the preceding paragraphs into a
single reference table that downstream reviewers can
consult when tempted to read a \texttt{MioCertificate} as
more than its schema permits.

\begin{center}
\small
\begin{tabular}{ll}
\hline\hline
Intended reading & Enforced status \\
\hline
``MIO certified the model''
  & prohibited (not a truth certificate) \\
``MIO evidence + HTT evidence''
  & prohibited (no master score) \\
``MIO posterior mean''
  & prohibited (no posterior field) \\
``HTT likelihood $\times$ MIO score''
  & prohibited (ingestion refused) \\
``MIO admissibility gate''
  & prohibited (TSC role) \\
\hline
``MIO measurement + caveats''
  & permitted (departure report) \\
``MIO $\leftrightarrow$ HTT cross-check report''
  & permitted (agreement only) \\
``MIO advisory flag on HTT workflow''
  & permitted (domain caveat) \\
\hline\hline
\end{tabular}
\end{center}

Each prohibited row corresponds to an explicit code-level
defence, as summarised in §12.6 and elaborated in
appendix~A40.  Each permitted row corresponds to a
currently wired cross-check channel
(appendix~A34.3).

\paragraph{Reviewer checklist.}

For a reviewer confronted with a MIO figure or a
\texttt{MioCertificate} JSON artefact, the following checks
summarise the present section:

\begin{itemize}
\item The \texttt{report\_type} names one of the five
  Phase-J reductions;
  \texttt{reduction\_status} is one of
  \{\texttt{theory-direct},
  \texttt{theory-approximate},
  \texttt{diagnostic-only}\}.
\item The filename begins with \texttt{mio\_} and the
  figure caption carries a
  ``[model-independent]'' tag (v3 §14.1).
\item No field on the certificate carries posterior,
  likelihood, or evidence semantics.
\item The certificate is not consumed as an HTT likelihood
  input anywhere in the downstream pipeline; any such
  consumption fails at the type-check layer.
\item The reader interprets the certificate as a
  data-derived measurement, not as a truth certificate on
  any cosmological hypothesis.
\end{itemize}

A figure or certificate that fails any of the five items is
either miscaptioned or indicates a G19 violation; either
case should prompt a correction before the figure is
referenced as evidence.
